% GENERATED FILE. Edit canonical lessons, then run npm run generate:books.
% canonical-source-hash: fnv1a64:5387f3842e850c1a
% canonical-lessons: ZH-C210-zise, ZH-C210-zongse, ZH-C210-kuanshi, ZH-C210-shiyijian, ZH-C210-xuan

\chapter{Colours and Styles}
\label{ch:zh-colours-and-styles}
\hlchaptermodality{210}

\begin{chapteropening}
\textbf{By the end of this chapter:} \emph{I can name purple and brown, talk about styles, choose one and find the fitting room.}

Complete the last lesson of chapter 210: I can name purple and brown, talk about styles, choose one and find the fitting room.
\end{chapteropening}

\section[zǐsè]{zǐsè — purple}

\label{lesson:ZH-C210-zise}



\begin{quote}
Before the new one: say the Mandarin for to disturb, then the Mandarin for to permit.
\end{quote}

\subsection*{You'll want to know: zǐsè}
\emph{zǐsè} — \textquotedblleft{}purple\textquotedblright{}.

\begin{quote}
\emph{wǒ xǐhuan zǐsè} — I like purple.
\end{quote}

Its characters come later; for now it is a word for the ear.

\begin{cousinweb}
\emph{zǐsè} is built from \emph{zǐ} 'purple' and \emph{sè} 'colour', both inherited from Middle Chinese, \emph{tsjeX} and \emph{srik} (Guangyun, in Baxter's transcription), which glosses the first as a colour between the pure ones.
\end{cousinweb}

\begin{grammarlens}[title={What you've built}]
One more word for asking for what I want.
\end{grammarlens}

\subsection*{Your turn}
\begin{itemize}
\raggedright
  \item \emph{Say it:} \emph{zǐsè}
  \item \emph{Say it:} \emph{zǐsè}, once more
  \item \emph{Say it:} \emph{wǒ xǐhuan zǐsè} — I like purple.
  \item \emph{Recall:} say the Mandarin for to disturb, then the Mandarin for to permit, then say \emph{zǐsè} again
\end{itemize}

\begin{tcolorbox}[breakable,colback=teal!4,colframe=teal!35!black,title={Before you move on}]
What does zǐsè mean? (\textquotedblleft{}Purple\textquotedblright{}.) Say it once more.
\end{tcolorbox}
\section[zōngsè]{zōngsè — brown}

\label{lesson:ZH-C210-zongse}



\begin{quote}
Before the new one: say the Mandarin for to permit, then the Mandarin for purple.
\end{quote}

\subsection*{You'll want to know: zōngsè}
\emph{zōngsè} — \textquotedblleft{}brown\textquotedblright{}.

\begin{quote}
\emph{wǒ yào zōngsè de} — I'll have the brown one.
\end{quote}

Its characters come later; for now it is a word for the ear.

\begin{cousinweb}
\emph{zōngsè} is built from \emph{zōng} and \emph{sè} 'colour'. The second is inherited from Middle Chinese \emph{srik} (Guangyun, in Baxter's transcription); no dictionary we use records an older form of the first, so its history is not settled.
\end{cousinweb}

\begin{grammarlens}[title={What you've built}]
One more word for asking for what I want.
\end{grammarlens}

\subsection*{Your turn}
\begin{itemize}
\raggedright
  \item \emph{Say it:} \emph{zōngsè}
  \item \emph{Say it:} \emph{zōngsè}, once more
  \item \emph{Say it:} \emph{wǒ yào zōngsè de} — I'll have the brown one.
  \item \emph{Recall:} say the Mandarin for to permit, then the Mandarin for purple, then say \emph{zōngsè} again
\end{itemize}

\begin{tcolorbox}[breakable,colback=teal!4,colframe=teal!35!black,title={Before you move on}]
What does zōngsè mean? (\textquotedblleft{}Brown\textquotedblright{}.) Say it once more.
\end{tcolorbox}
\section[kuǎnshì]{kuǎnshì — style; design}

\label{lesson:ZH-C210-kuanshi}



\begin{quote}
Before the new one: say the Mandarin for purple, then the Mandarin for brown.
\end{quote}

\subsection*{You'll want to know: kuǎnshì}
\emph{kuǎnshì} — \textquotedblleft{}style; design\textquotedblright{}.

\begin{quote}
\emph{zhège kuǎnshì hěn hǎokàn} — This style looks nice.
\end{quote}

Its characters come later; for now it is a word for the ear.

\begin{cousinweb}
\emph{kuǎnshì} is built from \emph{kuǎn} and \emph{shì} 'a model', both inherited from Middle Chinese, \emph{khwanX} and \emph{syik} (Guangyun, in Baxter's transcription), which glosses the first 'sincere' and the second 'a rule'; how the pair came to mean a style is not settled.
\end{cousinweb}

\begin{grammarlens}[title={What you've built}]
One more word for asking for what I want.
\end{grammarlens}

\subsection*{Your turn}
\begin{itemize}
\raggedright
  \item \emph{Say it:} \emph{kuǎnshì}
  \item \emph{Say it:} \emph{kuǎnshì}, once more
  \item \emph{Say it:} \emph{zhège kuǎnshì hěn hǎokàn} — This style looks nice.
  \item \emph{Recall:} say the Mandarin for purple, then the Mandarin for brown, then say \emph{kuǎnshì} again
\end{itemize}

\begin{tcolorbox}[breakable,colback=teal!4,colframe=teal!35!black,title={Before you move on}]
What does kuǎnshì mean? (\textquotedblleft{}Style; design\textquotedblright{}.) Say it once more.
\end{tcolorbox}
\section[\texorpdfstring{\zh{试衣间} (shìyījiān)}{试衣间 (shìyījiān)}]{\zh{试衣间} (shìyījiān) — fitting room}

\label{lesson:ZH-C210-shiyijian}



\begin{quote}
Before the new one: say the Mandarin for brown, then the Mandarin for style of clothes.
\end{quote}

\subsection*{You'll want to know: \zh{试衣间}}
\textbf{\zh{试衣间}} — \emph{shìyījiān} — \textquotedblleft{}fitting room\textquotedblright{}.

\begin{quote}
\textbf{\zh{我在找试衣间}} \emph{wǒ zài zhǎo shìyījiān} — I'm looking for the fitting room.
\end{quote}

\begin{cousinweb}
\textbf{\zh{试衣间}} (\emph{shìyījiān}) is built from \zh{试} 'to try', \zh{衣} 'clothes' and \zh{间} 'a room'. The first two are inherited from Middle Chinese, \emph{syiH} and \emph{'j+j} (Guangyun, in Baxter's transcription).
\end{cousinweb}

\begin{grammarlens}[title={What you've built}]
One more word for asking for what I want.
\end{grammarlens}

\subsection*{Your turn}
\begin{itemize}
\raggedright
  \item \emph{Say it:} \emph{shìyījiān}
  \item \emph{Say it:} \emph{shìyījiān}, once more
  \item \emph{Say it:} \emph{wǒ zài zhǎo shìyījiān} — I'm looking for the fitting room.
  \item \emph{Recall:} say the Mandarin for brown, then the Mandarin for style of clothes, then say \emph{shìyījiān} again
\end{itemize}

\begin{tcolorbox}[breakable,colback=teal!4,colframe=teal!35!black,title={Before you move on}]
What does \zh{试衣间} mean? (\textquotedblleft{}Fitting room\textquotedblright{}.) Say it once more.
\end{tcolorbox}
\section[xuǎn]{xuǎn — to choose}

\label{lesson:ZH-C210-xuan}



\begin{quote}
Before the new one: say the Mandarin for style of clothes, then the Mandarin for fitting room.
\end{quote}

\subsection*{You'll want to know: xuǎn}
\emph{xuǎn} — \textquotedblleft{}to choose\textquotedblright{}.

\begin{quote}
\emph{nǐ xuǎn nǎge yánsè} — Which colour are you choosing?
\end{quote}

Its characters come later; for now it is a word for the ear.

\begin{cousinweb}
\emph{xuǎn} is inherited from Middle Chinese \emph{sjwenX} (Guangyun, in Baxter's transcription), glossed 'to choose'.
\end{cousinweb}

\begin{grammarlens}[title={What you've built}]
One more word for asking for what I want.
\end{grammarlens}

\subsection*{Your turn}
\begin{itemize}
\raggedright
  \item \emph{Say it:} \emph{xuǎn}
  \item \emph{Say it:} \emph{xuǎn}, once more
  \item \emph{Say it:} \emph{nǐ xuǎn nǎge yánsè} — Which colour are you choosing?
  \item \emph{Recall:} say the Mandarin for style of clothes, then the Mandarin for fitting room, then say \emph{xuǎn} again
\end{itemize}

\begin{tcolorbox}[breakable,colback=teal!4,colframe=teal!35!black,title={Before you move on}]
What does xuǎn mean? (\textquotedblleft{}To choose\textquotedblright{}.) Say it once more.
\end{tcolorbox}
